%- What is important
%- what is done
%- what is wrong
%- how do we do solve it

Pharmacokinetics (PK) studies the dynamics of how the human body processes a drug.
%
In a typical early-development setting, each subject receives a single dose, and the goal is to characterize how systemic drug concentration unfolds over time.
%
This characterization requires not only representing temporal dynamics, but also capturing \textit{heterogeneity} because, even under identical doses, different individuals can exhibit markedly different profiles, due to factors such as physiology, disease state, and co-medications.

The key difficulty is that these dynamics must be inferred from sparse and noisy measurements collected across a \textit{study}: a cohort of individuals, each observed at an irregular, and often subject-specific set of time points. 
%
Operationally, one would like to (i) predict the concentration profile of a new individual, or (ii) forecast the unobserved future of a patient from only a few early measurements.
%
PK formalizes this setting by treating each measured concentration as a noisy sample of an underlying, patient-specific, and continuous-time \textit{concentration function} generated by a dynamical system~\citep{gibaldi1982pharmacokinetics}.
%
In the population setting, this is most often expressed through nonlinear mixed-effects (NLME) models, which decompose variability into fixed effects shared across the cohort, and random effects capturing individual deviations~\citep{lavielle2014}.
%
From this viewpoint, the fundamental inferential object is a \textit{distribution over plausible concentration functions}, conditioned on study information and (optionally) partial observations.

Recent work has begun to approximate such concentration functions with deep generative models, such as neural ODEs~\citep{lu2021deep}. 
%
However, PK datasets are not only sparse and noisy, but also small (especially in Phase I and many Phase II settings). 
%
Training a high-capacity model from scratch for each new dataset is therefore both statistically fragile and operationally expensive. 
%
To bridge the gap between expressive modeling and data scarcity, we turn to \textit{amortized inference}.

Amortization shifts the computational cost of inference from repeated, dataset-specific optimization to a single, up-front \textit{pretraining} phase across large and diverse datasets --- typically synthetically generated from \textit{mechanistic priors}. 
%
The heterogeneous pretraining distributions encourage models to learn \textit{reusable inference algorithms} that map a context (\textit{e.g.}, the study and partial observations) directly to a target (posterior) distribution.
%
This paradigm underlies prior-fitted networks (PFNs) for tabular data~\citep{muller2022transformers, hollmann2023tabpfn, hollmann2025accurate}, foundation inference models (FIMs) for system identification~\citep{berghaus2024foundation, fim_sde, berghaus2026incontext, hubers2026foundation}, and many foundation models for time-series imputation~\citep{seifner2025zeroshot} and forecasting~\citep{dooley2024forecastpfn, bhethanabhotla2024mamba4cast, hemmer2025true}.

\texttt{AICMET}~\citep{marin2025amortized} represents a first step in this direction for PK.
%
It was pretrained on synthetic studies generated from stochastic priors on compartmental models, and performs \textit{zero-shot posterior estimation} over concentration functions using a hierarchical mixed-effects representation.
%
The latter is implemented through neural-process-style latent variables~\citep{garnelo2018neural}.
%
A limitation of this design is the reliance on a compressed representation of the study context, which can lead to underfitting when the context is complex~\citep{kim2019attentive}, and to uncertainty that is not sharply calibrated~\citep{wang2025bridge}.

In this work, we dispense from explicit mixed-effect parametrizations and instead learn the distribution over concentration functions directly in function space.
%
We introduce \textit{Prior-fitted Functional Flows} (PFFs), which leverages conditional optimal-transport (OT) flow matching~\citep{kerrigan2024dynamic} to transport a simple reference measure over functions (a Gaussian process) into the conditional distribution over drug-concentration functions. 
%
The construction is conditional at two levels: first at the level of the study (\textit{i.e.}, the cohort-level context), and second, at the level of a patient’s early measurements for forecasting.
%
We therefore make the following contributions:

\begin{enumerate}
    \item We release an open-access reference dataset mined from clinical bioequivalence studies, containing inferred PK parameters and synthetic concentration--time samples used to calibrate physiologically plausible pretraining priors.

    \item We introduce \texttt{PFF}, a prior-fitted functional flow model for zero-shot PK inference. PFF directly transports measures over concentration--time functions, combines study-level conditioning with GP posterior references for partially observed individuals, and uses continuum-attention neural operators~\citep{calvello2024continuum}. Experiments and ablations show improved forecasting and sample fidelity over NLME, \texttt{AICMET}, and PFF variants.
\end{enumerate}

% In this work, we leverage the success of transformer architectures for in‑context Bayesian inference \citep{mittal2025amortized, muller2021transformers, berghaus2024foundation, seifner2025zeroshot, seifner2025foundation} and translate these methodologies to mixed‑effects pharmacokinetic models. We train a transformer decoder that operates in context: it can both generate synthetic patient trajectories and, crucially, exploit the collective context provided by previously profiled trial participants to deliver dose‑conditioned predictions for newly enrolled patients receiving the same drug. The ability to generate synthetic patients for specific populations on a given drug allows us to characterize the inter‑patient variability required to optimize individualized dosing regimens, evaluate safety margins under covariate shifts, and quantify the predictive uncertainty needed for adaptive trial design and regulatory submission.

% Similar to \citet{seifner2025zeroshot, seifner2025foundation} and \citet{berghaus2024foundation}, we first propose a statistical model that imposes stochastic differential equation priors on classical compartment ODE model parameters. Furthermore, we introduce a hidden‑variable formulation akin to neural processes \citep{garnelo2018neural, nguyen2022transformer}, but with a hierarchical structure: a global latent representation captures population‑level effects, while local latents model individual‑specific deviations. Our transformer based decoder integrates neural operator ideas \citep{lu2021learning,seifner2025zeroshot,seifner2025foundation} by encoding continuous time directly in the attention queries.

% By pre‑training on a massive corpus of simulated PK trajectories, the model internalizes model-specific pharmacological inductive biases. Consequently, bringing a new compound online involves only a brief, sometimes zero‑shot, adaptation step---shrinking model development timelines from weeks to seconds and freeing domain experts to influence therapeutic decisions rather than hand‑crafting ODE systems.
% In short, we provide:  (i) a unified probabilistic framework combining compartmental PK models with hierarchical latent-variable transformers to capture both population- and subject-level structures; (ii) an amortized in-context inference mechanism enabling calibrated predictions from as little as a single observation per individual; (iii) a dose-aware neural architecture that models continuous-time dynamics and naturally handles irregular sampling; and (iv) a pretrained generative model supporting zero-shot forecasting, safety margin estimation, and predictive uncertainty quantification.
%
% Together, these contributions outline a data‑efficient, population‑aware, and mechanistically grounded pathway towards truly personalized pharmacotherapy.
